% capitulo1.tex
\chapter{Introducción}
\section{Antecedentes}
El diseño de las placas base en estructuras de acero es un elemento crítico para la transferencia de cargas de la superestructura a la cimentación. Cambios en las disposiciones normativas y en los criterios de diseño (resistencia, ductilidad y verificación sísmica) pueden afectar dimensiones de placa, disposición y resistencia de pernos de anclaje, uso de rigidizadores y detalles constructivos.

En esta tesina se realiza una comparativa entre la NTC 2017 y la NTC 2023 aplicada al diseño de placas base para estructuras de acero, evaluando diferencias en hipótesis de carga, factores de reducción, prescripciones para anclajes y criterios de verificación. A continuación se muestran dos figuras que ayudan a contextualizar el problema y su origen sísmico.

\begin{figure}[H]
  \centering
  \includegraphics[width=0.4\textwidth]{\detokenize{figuras/001_Vista en corte de una conexión columna-cimentación columna de acero, placa base, pernos de anclaje, grout y pedestal de concreto..png}}
  \caption{Vista en corte de una conexión columna-cimentación columna de acero, placa base, pernos de anclaje, grout y pedestal de concreto.}
  \label{fig:baseplate}
\end{figure}

\begin{figure}[H]
  \centering
  \includegraphics[width=0.4\textwidth]{\detokenize{figuras/002_Mapa de la República Mexicana señalando el epicentro del sismo del 19-S 2017 y la Ciudad de México..png}}
  \caption{Mapa de la República Mexicana señalando el epicentro del sismo del 19-S 2017 y la Ciudad de México.}
  \label{fig:map2017}
\end{figure}

Las Figuras~\ref{fig:baseplate} y~\ref{fig:map2017} ilustran, respectivamente, el detalle constructivo objeto de estudio (placa base y pernos de anclaje) y el evento sísmico que motivó revisiones normativas posteriores que pueden influir en criterios de diseño.
\section{Objetivos}
\subsection{Objetivo general}
Enunciar el objetivo general.
\subsection{Objetivos específicos}
\begin{itemize}
  \item Objetivo específico 1.
  \item Objetivo específico 2.
\end{itemize}
\section{Justificación}
Justificación de la investigación y su aporte a la ingeniería civil.
